\paragraph{Lifting an outer coordinate through a quadratic component.}
\label{par:row-mixed}
\hypertarget{source-row-mixed}{}

\begin{lemma}[Lifting the main tensors]\label{lem:mixed-row-tensors}
For the main tensors on a one-sign component, or on a mixed
component with \(e=2\), put
\(y_c=\lceil\eta+1/2+\log_pH\rceil\).
For each such tensor, let \(y\in N_c(T_2)\) represent its image
under the forward comparison \(\alpha\) of Section~\ref{par:row-metric},
modulo \(XN_c(T_2)\). There exists \(v\in N_c(T_1)\) with
\[
 (\gamma-1)v\in XN_c(T_1),\qquad Fv\equiv p^{\rho+\max(2\sigma,y_c)}y\pmod{XN_c(T_2)}.
\]
\end{lemma}

\begin{proof}
Let \(f\otimes x\) be one of the simple tensors above. We require an
invariant in the source with the same image as a power of \(p\) times
its image under \(\alpha\) modulo \(XN_c(T_2)\). Two operations suffice.
First, if an invariant model tensor has an invariant source partner
under \(F\), apply \(\alpha\) and use naturality. Changing the partner by \(u\in X^M E^{3H}P_c\) preserves its cut class because
\[
 X^M E^{3H}P_c\subseteq XN_c.
\]

Second, if a model
tensor belongs to \(N^{\rm mod}_2\), its image under \(\alpha\) belongs to
\(N_c(T_2)\), hence has an integral lift after multiplication by
\(p^\rho\). Any such lift is invariant modulo \(XN_c(T_1)\).
Here a model tensor with a scalar coefficient and a preceding constant
can be lifted by multiplying the scalar by a lift of that constant,
provided the resulting tensor meets the invariance test. Actual
preceding constants have lifts after \(p^\rho\) by the same Wach
cokernel bound on the preceding component.

\proofstep{Components with one sign.}
At a component with only plus tests, the product of each
coefficient of \(f\) with \(x\) lies in the required
filtration, since \(f\) starts in degree \(-b_\tau\) or
later and \(x\in\operatorname{Fil}^{b_\tau}\). Thus
\[
 f\otimes x\in N^{\rm mod}_2.
\]

At a component with only minus tests, choose a preceding constant
\(x_1\) with
\[
 F x_1=p^\rho x,
\]
and take \(f\otimes x_1\).
The source conditions are
\[
  E^H f\in A,\qquad
  (\gamma-1)f\in XE^{-H}A,\qquad
  \operatorname{ord}_{\tau}\big((\gamma-1)f/X\big)
      \ge b_\tau+S .
\]
The last inequality suffices because the entire source constant
space is in \(\operatorname{Fil}^{-b_\tau-S}\) at each minus test.

\proofstep{The two roots of the quadratic component.}
For opposite signs assume \(e=2\). Write \(r_U,r_D\) for the
two parameter values at the plus and minus tests, and put
\[
 y_U=X-r_U,\ y_D=X-r_D.
\]
Division by \(X-r\) at any such
value is integral since \(r\) belongs to the maximal ideal and
\(O\) is complete. Over the fraction field the changes to \(y_U\) and \(y_D\) have invertible linear term at their respective roots. Each preserves the order of vanishing and transports the filtration test. In particular equal \(X\)-values impose identical scalar Taylor conditions.

If a plus test has the minus
root as well, the high vanishing of \((\gamma-1)f/X\) at the
minus root also suffices at that test, since its positive degree
is greater than \(\eta\). This gives the preceding-lift argument
when every plus test shares a minus root. Otherwise there are just
two distinct roots, the minus tests at \(r_D\) and the remaining
plus tests at \(r_U\), with the same labeled filtration at all
points over the corresponding embedding.

Write
\[
 b=b_U,\qquad B=b_D+S,\qquad
 E=a_2 y_U y_D,\qquad v(a_2)=0,\qquad d_*=r_D-r_U.
\]
Both roots have valuation \(s/(p-1)=1/2\). Also
\(v(d_*)=1/2\): at \(r_U\) differentiate the cyclotomic identity
\(Q(\pi)={\rm unit}\cdot E(X(\pi))\), where
\(Q=((1+\pi)^p-1)/\pi\). The derivative of \(Q\) at a primitive
point has valuation \(1-1/(p-1)\), and \(X'(\pi)\) has valuation
\((s-1)/(p-1)\). Hence
\[
 v(E'(r_U))=v(Q'(\pi))-v(X'(\pi))=1-\frac1{p-1}-\frac{s-1}{p-1}=\frac12.
\]
This accounts for the half power in the following allowance.
Let \(y_c=\lceil\eta+1/2+\log_p H\rceil\).

\proofstep{Nearby breaks.}
Suppose first \(b\le B+2\eta\). The numerator \(a=E^H f\) is
divisible by \(y_U^{H-b}\) in \(A\). At the down root the series
of \(f\) equals a constant \(x_0\) modulo degree \(B\).
Consequently the coefficients of \(a/y_U^{H-b}\) there, through
degree \(H+B-1\), agree with those of
\(x_0 a_2^H y_D^H(y_D+d_*)^b\). They are integral in \(y_D\)
by division in \(A\). In particular
\[
           x_0a_2^H \binom{b}{j} d_*^{\,b-j}\in O
                   \quad (0\le j<B,\ j\le b).
\]

If \(b<B\) take \(j=b\). Otherwise take \(j=B-1\); then
\(b-j\le2\eta+1\) and
\[
 v_p\binom{b}{j}\le\lfloor\log_p b\rfloor
\]
by the factorial formula, a difference of at most one in each
power-\(p\) count. In either case \(p^{y_c}x_0\in O\).
Subtract this constant from \(p^{y_c}f\). The remainder times
\(x\) belongs to \(N^{\rm mod}_2\). At the minus points its
order is at least \(B\); at the plus points its order is at
least \(-b\) and the target constant is in \({\rm Fil}^b\).
The constant term itself lifts by an arbitrary preceding lift,
since a constant is fixed by \(\gamma\). Applying the two operations
gives an invariant lift in the source Wach lattice at exponent \(\rho+y_c\).

\proofstep{Separated breaks.}
Suppose now \(B+2\eta<b\). A division by \(d_*^{b-B+1}\) could
cost an amount proportional to the outer height. Instead form the
auxiliary numerator
\[
 a_0=y_U^{H-b}y_D^{H+B}.
\]
The tensor
\(a_0 x/E^H\) belongs to \(N^{\rm mod}_2\) by its orders at both
roots. Apply \(\alpha\), lift its image in the target Wach lattice after \(p^\rho\)
in \(N_c(T_1)\), and apply \(\beta\).

Write the result with denominator
\(E^H\) as \(\sum_j a_j v_j/E^H\), \(a_j\in A\), where \(v_j\)
are integral preceding constants of the source. The sum in every
numerator coefficient of degree \(i\) at a root lies in
\({\rm Fil}^{H-i}_\tau\). On applying \(F\) its expansion in degrees
through \(H\) agrees with
\[
                  p^{\rho+2\sigma}a_0 x .
\]
Indeed the forward and backward composites give this scalar and the
coordinate modulus, after multiplying by \(E^H\), contains the
factor \(E^{4H}\) in every error term. Integral bases on the preceding
component give a representation in this common constant space for
both filtrations.

Expand each \(a_j\) at \(r_D\) through degrees \(<H-\eta\), with
coefficients in \(O\), and remove this polynomial in \(y_D\).
In degree \(i<H-\eta\) the corresponding sum against \(v_j\)
belongs to
\[
 {\rm Fil}^{H-i}_D=0,
\]
since the largest source
jump there is at most \(\eta\). These polynomial coefficients
therefore give zero sums against both \(v_j\) and \(Fv_j\) at the
plus root also, by the common crystalline constant identification.
Divide the remaining functions by \(y_D^{H-\eta}\), obtaining
integral \(a'_j\). At \(r_U\) this divisor is a unit over the
fraction field and the source coefficient condition
\({\rm Fil}^{H-i}_U\) persists. The new target expansion agrees
through degree \(H\) with
\[
       p^{\rho+2\sigma}y_U^{H-b}
                  (y_U+r_U-r_D)^{B+\eta}x .
\]

Take its coefficient of degree \(j_*=H-b+B+\eta\) in \(y_U\).
The assumption gives \(j_*<H-\eta\); changing from \(t\) to \(y_U\)
preserves the filtration test and precision, by triangular substitution.
This is the leading coefficient of the displayed polynomial. Hence
the same integral linear combination of the \(v_j\) maps to
\(p^{\rho+2\sigma}x\) and belongs to
\({\rm Fil}^{H-j_*}_U\subset{\rm Fil}^{\eta+1}_U
={\rm Fil}^b_U\). The equality uses the empty interval of jumps
strictly between \(\eta\) and \(b\). This combination supplies a
preceding lift compatible with the plus test. Keeping \(f\) gives
an invariant model lift at both signs, hence an invariant Wach-lattice
lift after \(\alpha\). Altogether a single exponent
\[
                       y_1=\rho+\max(2\sigma,y_c)
\]
serves every main tensor on every component, including the one-sign
cases.  \label{cor:row-mixed-1}
\end{proof}
